\chapter{Канонический двойственный каркас и совместные моменты одной нити}
\label{ch:v6-single-thread-connectivity-weighted-moment-parent-gate}

\section{Постановка}

Предыдущая глава показала, что фазы остаточной $Z_3$-голономии не могут
компенсировать неравные плотности четырёх локальных проходов. Это закрывает
сырой и чисто фазовый третий момент, но не связностный: ориентация входящего
прохода и правило его продолжения через локальный узел образуют пару
первичного и двойственного направлений.

Канонический двойственный каркас не является новым физическим полем. В
теории каркасов он однозначно строится обращением оператора каркаса
\cite{single-thread-duffin-schaeffer1952}. В проекте необходимый оператор
уже присутствует как локальный второй момент $R$.

\section{Тетраэдральный каркас из существующей коизометрии}

Существующая аффинная коизометрия $V:\mathbb R^4\to\mathbb R^3$ имеет
свойства
\begin{equation}
 VV^{\mathsf T}=I_3,
 \qquad
 V^{\mathsf T}V=P_3=I_4-\frac14\mathbf1\mathbf1^{\mathsf T}.
\end{equation}
После нормировки четыре столбца $V$ дают направления $n_a$ с матрицей
Грама
\begin{equation}
 n_a\cdot n_b=
 \begin{cases}
 1,&a=b,\\
 -1/3,&a\ne b.
 \end{cases}
\end{equation}
Следовательно, это в точности правильный тетраэдральный каркас, уже
содержащийся в родителе $4=1+3$. Отдельная четвёрка направлений не вводится.

\section{Первичный и двойственный проходы}

Пусть положительные локальные доли проходов равны $w_a$ и
\begin{equation}
 R=\sum_{a=1}^{4}w_a n_an_a^{\mathsf T}.
\end{equation}
Это тот же второй момент, бесследовая часть которого задаёт $Q$. Пока все
четыре веса положительны, $R$ обратим. Канонические двойственные направления
равны
\begin{equation}
 d_a=R^{-1}n_a.
\end{equation}
Они удовлетворяют формуле восстановления
\begin{equation}
 \sum_aw_a d_an_a^{\mathsf T}=I_3.
\end{equation}

Связностный третий момент должен учитывать не три независимых измерения
одного и того же времени пребывания, а один двойственный индекс локальной
сшивки и два первичных ориентационных индекса:
\begin{equation}
 \mathcal M_3
 =\sum_aw_a\operatorname{Sym}
 \bigl(d_a\otimes n_a\otimes n_a\bigr).
\end{equation}

\begin{theorem}[Точное восстановление тетраэдрического момента]
Для рабочей $C_3$-упорядоченной семьи
\begin{equation}
 w=(p,q,q,q),
 \qquad q=\frac{1-p}{3},
 \qquad 0<p<1,
\end{equation}
выполняется тождество
\begin{equation}
 \operatorname{STF}\mathcal M_3
 =\frac34\sum_{a=1}^{4}n_a^{\otimes3}
 =\frac34T_*.
\end{equation}
Коэффициент $3/4$ не зависит от $p$ и не является подгоночным параметром.
\end{theorem}

\begin{proof}
Тетраэдральные тождества дают
$\sum_an_a=0$ и $\sum_an_an_a^{\mathsf T}=4I_3/3$. Поэтому $R$ имеет одно
собственное значение вдоль выбранной оси и двукратное поперечное
собственное значение. Подстановка соответствующего $R^{-1}$ в
$\mathcal M_3$ показывает, что вся зависимость от $p$ остаётся только в
следовых частях ранга три. Проекция $\operatorname{STF}$ их удаляет, а
оставшаяся бесследовая часть равна $3T_*/4$. Компонентная символическая
редукция даёт нулевой остаток тождества.
\end{proof}

Для рабочего спектра
\begin{equation}
 w=(0.9011874254,0.03293752488^{(3)})
\end{equation}
численный остаток равен $3.68\cdot10^{-16}$, тогда как сырой момент давал
$0.6133541143$, а лучший нетривиальный фазовый подъём --- $0.6671606493$.
Скан $256$ случайных положительных четвёрок сохранил коэффициент $3/4$ с
максимальным остатком $1.23\cdot10^{-15}$. Скан $1001$ точки всей рабочей
$C_3$-семьи дал остаток не выше $3.48\cdot10^{-13}$ даже у почти
вырожденной границы.

\section{Почему это не запрещённый обратный вес}

Ранее $R^{-1}$ был отвергнут как автоматически добавляемый энергетический
барьер: родительское действие не выводило такую энергию. Здесь обратный
оператор играет другую роль. Он является единственным каноническим
двойственным каркасом для уже заданных локальных проходов и входит в
формулу восстановления связности, а не в потенциал или свободную энергию.

Тем самым отношение $p/q=27.36050837$ не вводится вручную. Оно поглощается
полным матричным оператором $R^{-1}$, а бесследовый ответ после
симметризации вообще не зависит от этого отношения.

\begin{nogocriterion}
Полученное локальное тождество нельзя объявлять доказательством одной
глобальной нити. Оно показывает совместимость локальных полей $Q$ и $T$ с
одним набором проходов, но ещё не сшивает эти проходы в единственный
замкнутый цикл и не выводит конечную толщину или планковский масштаб.
\end{nogocriterion}

\section{Следующий различающий тест}

Локальный моментный запрет снят. Следующий гейт должен построить глобальное
правило сшивки: ориентированный граф переходов между тетраэдральными
проходами обязан иметь один эйлеров цикл, а не четыре независимые нити или
несколько несвязных циклов.

Нулевая гипотеза: существующие аффинная частичная изометрия, Real-пара и
$Z_3$-голономия однозначно задают сбалансированный связный поток, допускающий
единственный глобальный цикл.

Стоп-критерий: если локальные степени можно сбалансировать только ценой
нескольких компонент связности или нового непрерывного переходного веса,
модель одной глобальной нити не выведена.

\begin{protocol}
\begin{enumerate}
 \item аффинная коизометрия задаёт правильный тетраэдр: \textbf{да};
 \item оператор второго момента является оператором каркаса: \textbf{да};
 \item канонический двойственный каркас существует: \textbf{при $w_a>0$};
 \item смешанный первично-двойственный момент даёт $T_*$: \textbf{точно};
 \item фиксированный коэффициент: \textbf{$3/4$};
 \item новый относительный параметр: \textbf{нет};
 \item $R^{-1}$ добавлен как энергия: \textbf{нет};
 \item локальный мост $Q+T$ одной нити: \textbf{пройден};
 \item единый глобальный цикл: \textbf{ещё не выведен};
 \item следующая проверка: \textbf{глобальная циклическая сшивка};
 \item рождение материи: \textbf{не закрыто}.
\end{enumerate}
\end{protocol}

Машинный сертификат сохранён в
\path{s2t_v6_single_thread_connectivity_weighted_moment_parent_gate_results.json}.

\begin{thebibliography}{9}
\bibitem{single-thread-duffin-schaeffer1952}
R.~J.~Duffin, A.~C.~Schaeffer,
\emph{A Class of Nonharmonic Fourier Series},
Transactions of the American Mathematical Society 72 (1952), 341--366.
\end{thebibliography}